% ============================================================
% Chapter 1: Introduction
%   Section: Research Questions
% Included from chapters/01-introduction/introduction.tex
% ============================================================

% Fixed by the author on 2026-10-14, replacing the roadmap's proposal.
% The author's own wording, kept so the formal versions below can be
% checked against it:
%   1. Can I create a retriever that uses the internet and simplifies
%      a journalist's finding tasks according to certain topics?
%   2. Can I verify the information with my own logic, and also verify
%      the news I retrieve with the structure we are creating - and how
%      can we quantify it?
%   3. How close can a person and an LLM get after getting the ...
%      (the sentence was cut off; RQ3 below completes it as "the same
%      news and evidence" - to be confirmed by the author)
%   4. All of this joined, a system which would in some way replace
%      the journalist's job?
% Each is answered by the metrics of Section~\ref{sec:rq-metrics};
% whatever answers none of them is cut or moved to future work.

\section{Research Questions}
\label{sec:research-questions}

The work is organised around four research questions. The first three
follow the journalist's own sequence --- finding the news, checking
it, and deciding --- and the fourth asks what the three amount to
when they are joined into one system.

\begin{description}[leftmargin=1.2cm,labelwidth=1cm]
    \item[RQ1] \textbf{Finding.} Can a retriever that reads the open
    web, scoped to a chosen set of topics, simplify a journalist's
    task of finding constructive news on those topics?
    \item[RQ2] \textbf{Verifying.} Can the claims of the news it
    retrieves be verified automatically by the pipeline proposed in
    this work --- claim selection, evidence retrieval on the open web
    and a verdict from a language model --- and how can the quality
    of that verification be quantified?
    \item[RQ3] \textbf{Person and model.} Given the same news and the
    same evidence, how close does a large language model come to a
    person's judgement, both when choosing which news is worth
    publishing and when deciding whether a claim is supported?
    \item[RQ4] \textbf{The system.} Joined into one system, to what
    extent can these components take over the work of a journalist
    --- finding, selecting and verifying constructive news --- and
    at which points is a person still needed?
\end{description}

RQ4 is deliberately asked as a question of extent rather than as a
yes-or-no question: whether a system ``replaces'' a journalist is not
measurable, whereas the share of the work it performs at a stated
quality, and the points at which it does not, are.
